\section{THUẬT TOÁN K-MEANS}

\subsection{Khái niệm / Ý tưởng}
K-means là một thuật toán phân cụm không giám sát dùng để chia tập dữ liệu thành $k$ cụm sao cho: các điểm trong cùng cụm thì giống nhau (gần tâm cụm) và các cụm khác nhau thì khác biệt nhau (tâm cụm xa nhau).

Thuật toán tìm ra $k$ ``tâm cụm'' (centroid) và lặp đi lặp lại 2 bước:
\begin{enumerate}
    \item Gán mỗi điểm vào cụm có tâm gần nhất (thường dùng khoảng cách Euclid).
    \item Cập nhật tâm cụm bằng giá trị trung bình của các điểm trong cụm đó.
\end{enumerate}

\subsection{Quy trình thực hiện phân cụm}
\begin{enumerate}
    \item \textbf{Chọn số cụm $k$:} Dựa vào kinh nghiệm hoặc phương pháp Elbow, Silhouette...
    \item \textbf{Khởi tạo $k$ tâm cụm ban đầu:} Chọn ngẫu nhiên $k$ điểm từ dữ liệu.
    \item \textbf{Bước gán cụm (Assignment Step):} Với mỗi điểm $x_i$, tính khoảng cách đến từng tâm $\mu_j$ và gán $x_i$ vào cụm có tâm gần nhất.
    \item \textbf{Bước cập nhật tâm (Update Step):} Với mỗi cụm $C_j$, cập nhật tâm:
    $$\mu_j = \frac{1}{|C_j|} \sum_{x_i \in C_j} x_i$$
    \item \textbf{Kiểm tra hội tụ:} Nếu tâm cụm không đổi hoặc đạt số vòng lặp tối đa thì dừng lại.
\end{enumerate}

\subsection{Mã giả (Pseudo-code K-means)}
\begin{lstlisting}[caption={Mã giả thuật toán K-means}]
Input: D (du lieu), k (so cum)
Output: Cac cum C1...Ck va tam mu1...muk

KMEANS(D, k):
1. Chon ngau nhien k diem lam tam ban dau mu_1...mu_k
2. Lap:
       Gan cum: Voi moi diem x trong D, gan vao cum Cj co dist(x, mu_j) nho nhat
       Cap nhat tam: mu_j = trung binh cac diem trong cum Cj.
   Cho den khi: Cac tam khong doi hoac dat so vong lap toi da.
3. Tra ve ket qua.
\end{lstlisting}

\subsection{Ưu / Nhược điểm của thuật toán K-means}
\textbf{Ưu điểm:}
\begin{itemize}
    \item \textbf{Đơn giản, dễ cài đặt:} Thuật toán chỉ dùng khoảng cách và giá trị trung bình.
    \item \textbf{Tốc độ nhanh, mở rộng tốt:} Độ phức tạp gần $O(k \times n \times t)$ (với $k$ là số cụm, $n$ số điểm, $t$ số vòng lặp).
    \item \textbf{Áp dụng tốt cho dữ liệu lớn:} Hoạt động rất tốt nếu dữ liệu có dạng cụm cầu (spherical clusters).
    \item \textbf{Kết quả trực quan:} Mỗi cụm có tâm giúp dễ diễn giải và làm tiền đề cho các bước phân tích sau.
\end{itemize}

\textbf{Nhược điểm:}
\begin{itemize}
    \item \textbf{Phải chọn $k$ trước:} Phải thử nhiều giá trị $k$ thông qua phương pháp Elbow hoặc Silhouette.
    \item \textbf{Nhạy với khởi tạo:} Tâm ban đầu khác nhau dẫn đến kết quả khác nhau, dễ rơi vào cực tiểu cục bộ.
    \item \textbf{Hạn chế về hình dạng cụm:} K-means phân cụm kém nếu cụm kéo dài hoặc có mật độ khác nhau.
    \item \textbf{Nhạy với nhiễu và ngoại lai:} Điểm outlier có thể kéo tâm cụm đi xa.
    \item \textbf{Dữ liệu số:} Chủ yếu dùng cho dữ liệu số, dữ liệu phân loại phải mã hóa phức tạp.
\end{itemize}